\chapter{Games, order, and surreal representatives}

\section{Finite games and birthdays}

\begin{definition}[Short combinatorial game]
  \label{def:game}
  \lean{Game,Game.left,Game.right,Game.zero,Game.one}
  \leanok
  A game is a finite rooted tree written \(x=\{X^L\mid X^R\}\), where the left- and
  right-option collections \(X^L\) and \(X^R\) are finite lists of games.  The basic games are
  \[
    0=\{\,\mid\,\},\qquad 1=\{0\mid\,\}.
  \]
\end{definition}

\begin{definition}[Birthday rank]
  \label{def:birthday}
  \lean{Game.birthday}
  \leanok
  \uses{def:game}
  The rank \(\bd(x)\in\mathbb N\) is defined recursively by
  \[
    \bd\bigl(\{X^L\mid X^R\}\bigr)
      = 1+\max\bigl(\{\bd(u):u\in X^L\}\cup
                         \{\bd(v):v\in X^R\}\cup\{0\}\bigr).
  \]
  Thus our convention is \(\bd(0)=1\).
\end{definition}

\begin{lemma}[Options have smaller birthday]
  \label{lem:birthday-options}
  \lean{Game.birthday_lt_of_mem_options,Game.birthday_lt_left,Game.birthday_lt_right}
  \leanok
  If \(u\in X^L\) or \(u\in X^R\), then \(\bd(u)<\bd(x)\).
\end{lemma}
\begin{proof}
  \leanok
  \uses{def:birthday}
  Write \(x=\{X^L\mid X^R\}\), and let \(B\) be the finite collection of
  birthdays of all members of \(X^L\) and \(X^R\).  If \(u\) belongs to either
  option collection, then \(\bd(u)\in B\); in particular, \(B\) is nonempty
  and
  \[
    \bd(u)\leq \max B.
  \]
  By Definition~\ref{def:birthday}, \(\bd(x)=1+\max B\), where adjoining
  \(0\) makes the same formula valid when there are no options.  Hence
  \(\bd(u)<\bd(x)\).  The left- and right-option assertions are the two
  specializations of this argument.
\end{proof}

\section{The recursive order}

\begin{definition}[Game comparison and equivalence]
  \label{def:game-order}
  \lean{Game.le,Game.lt,Game.eq}
  \leanok
  \uses{def:game,def:birthday,lem:birthday-options}
  The relation \(x\GameLe y\) is defined recursively by
  \[
    x\GameLe y
    \quad\Longleftrightarrow\quad
    \bigl(\forall x^L\in X^L,\ \neg(y\GameLe x^L)\bigr)
    \ \wedge\
    \bigl(\forall y^R\in Y^R,\ \neg(y^R\GameLe x)\bigr).
  \]
  Strict comparison and game equivalence are
  \[
    x\GameLt y \Longleftrightarrow x\GameLe y\ \wedge\ \neg(y\GameLe x),
    \qquad
    x\GameEq y \Longleftrightarrow x\GameLe y\ \wedge\ y\GameLe x.
  \]
\end{definition}

\begin{definition}[Well-founded induction schemes]
  \label{def:game-inductions}
  \lean{Game.R,Game.B,Game.T,Game.Q,Game.wf_R,Game.wf_B,Game.wf_T,Game.wf_Q}
  \leanok
  \uses{def:birthday}
  The development uses well-founded relations on one, two, three, and four games.  Their
  measures are respectively a birthday and the sums of two, three, and four birthdays.
\end{definition}

\begin{theorem}[Order calculus for games]
  \label{thm:game-order-calculus}
  \lean{Game.le_congr,Game.eq_congr,Game.le_trans',Game.le_trans,Game.eq_symm,Game.eq_trans,Game.lt_trans,Game.lt_of_lt_of_le,Game.lt_of_le_of_lt}
  \leanok
  The relation \(\GameLe\) is reflexive and transitive.  Consequently \(\GameEq\) is an
  equivalence relation, \(\GameLt\) is transitive, and strict and weak inequalities compose in
  either order.
\end{theorem}
\begin{proof}
  \leanok
  \uses{lem:birthday-options,def:game-order,def:game-inductions}
  For reflexivity, use the one-game induction scheme of
  Definition~\ref{def:game-inductions}.  Suppose reflexivity is known for all
  games of birthday smaller than \(\bd(x)\), and unfold the comparison in
  Definition~\ref{def:game-order}.  If \(x^L\) is a left option and
  \(x\GameLe x^L\), then the left-option clause of this last comparison says
  \(\neg(x^L\GameLe x^L)\).  On the other hand,
  Lemma~\ref{lem:birthday-options} gives \(\bd(x^L)<\bd(x)\), so the induction
  hypothesis gives \(x^L\GameLe x^L\), a contradiction.  If \(x^R\) is a
  right option and \(x^R\GameLe x\), the right-option clause similarly says
  \(\neg(x^R\GameLe x^R)\), again contradicting the induction hypothesis.
  Thus \(x\GameLe x\).

  For transitivity, use the three-game induction scheme of
  Definition~\ref{def:game-inductions}.  Assume \(a\GameLe b\) and
  \(b\GameLe c\), and induct on
  \(\bd(a)+\bd(b)+\bd(c)\).  If \(a^L\) is a left option and
  \(c\GameLe a^L\), then the induction hypothesis applied to the triple
  \((b,c,a^L)\) gives \(b\GameLe a^L\).  This contradicts the left-option
  clause of \(a\GameLe b\) in Definition~\ref{def:game-order}.  Similarly, if
  \(c^R\) is a right option and \(c^R\GameLe a\), the induction hypothesis for
  \((c^R,a,b)\) gives \(c^R\GameLe b\), contradicting the right-option clause
  of \(b\GameLe c\).  Both recursive calls decrease because
  Lemma~\ref{lem:birthday-options} replaces one coordinate by an option of
  smaller birthday.  Hence \(a\GameLe c\).

  The remaining assertions follow directly from the definitions in
  Definition~\ref{def:game-order}.  Symmetry and transitivity of \(\GameEq\)
  are obtained by composing the two weak inequalities in the appropriate
  directions.  For strict transitivity, compose the forward weak inequalities.
  If the reverse weak inequality held, transitivity would contradict one of
  the two strict hypotheses.  The mixed strict--weak laws follow from the same
  argument.
\end{proof}

\begin{theorem}[Equivalence from equivalent option lists]
  \label{thm:option-extensionality}
  \lean{Game.eq_of_equiv_options}
  \leanok
  Let \(a,b\) be games.  Suppose every left option of either game is equivalent to a left option
  of the other, and likewise every right option is equivalent to a right option of the other.
  Then \(a\GameEq b\).
\end{theorem}
\begin{proof}
  \leanok
  \uses{def:game-order,thm:game-order-calculus}
  We first verify the two clauses of Definition~\ref{def:game-order} for
  \(a\GameLe b\).  Let \(a^L\) be a left option of \(a\), and suppose that
  \(b\GameLe a^L\).  Choose a left option \(b^L\) of \(b\) such that
  \(a^L\GameEq b^L\).  The forward inequality in this equivalence and
  transitivity from Theorem~\ref{thm:game-order-calculus} give
  \(b\GameLe b^L\).  This is impossible: reflexivity from the same theorem
  gives \(b\GameLe b\), whose left-option clause forbids
  \(b\GameLe b^L\).

  Next let \(b^R\) be a right option of \(b\), and suppose
  \(b^R\GameLe a\).  Choose a right option \(a^R\) of \(a\) such that
  \(b^R\GameEq a^R\).  The reverse inequality \(a^R\GameLe b^R\) in this
  equivalence and transitivity
  give \(a^R\GameLe a\), contrary to the right-option clause of the reflexive
  comparison \(a\GameLe a\).  Thus \(a\GameLe b\).  Interchanging \(a\) and
  \(b\) proves \(b\GameLe a\), and Definition~\ref{def:game-order} then gives
  \(a\GameEq b\).
\end{proof}

\section{Surreal games}

\begin{definition}[Surreal game]
  \label{def:is-surreal}
  \lean{IsSurreal,IsSurreal.isSurreal_left,IsSurreal.isSurreal_right,IsSurreal.isSurreal_zero,IsSurreal.isSurreal_one}
  \leanok
  \uses{def:game,def:birthday,lem:birthday-options,def:game-order}
  A game \(x=\{X^L\mid X^R\}\) is surreal when every option is surreal and no right option is
  weakly below a left option:
  \[
    \forall x^L\in X^L\ \forall x^R\in X^R,
      \quad \neg(x^R\GameLe x^L).
  \]
  The games \(0\) and \(1\) are surreal, and every option of a surreal game is surreal.
\end{definition}

\begin{theorem}[A surreal game lies between its options]
  \label{thm:between-options}
  \lean{IsSurreal.left_lt,IsSurreal.lt_right,IsSurreal.xL_x_xR}
  \leanok
  If \(x\) is surreal, then
  \[
    \forall x^L\in X^L,\quad x^L\GameLt x,
    \qquad
    \forall x^R\in X^R,\quad x\GameLt x^R.
  \]
\end{theorem}
\begin{proof}
  \leanok
  \uses{lem:birthday-options,def:game-order,def:game-inductions,thm:game-order-calculus,def:is-surreal}
  We prove the assertion for left options; the right-option proof is its
  mirror.  Use the one-game induction scheme of
  Definition~\ref{def:game-inductions}.  Let \(x\) be surreal, and let
  \(x^L\) be a left option.  By Definition~\ref{def:is-surreal}, \(x^L\) is
  itself surreal.  We verify the two clauses of
  Definition~\ref{def:game-order} for \(x^L\GameLe x\).

  Let \(z\) be a left option of \(x^L\), and suppose \(x\GameLe z\).  Since
  \(\bd(x^L)<\bd(x)\) by Lemma~\ref{lem:birthday-options}, the induction
  hypothesis applied to the surreal game \(x^L\) gives \(z\GameLe x^L\).
  But the left-option clause of \(x\GameLe z\), applied to the option \(x^L\)
  of \(x\), forbids \(z\GameLe x^L\).  Next let \(x^R\) be a right option of
  \(x\).  The separation condition in Definition~\ref{def:is-surreal} is
  precisely \(\neg(x^R\GameLe x^L)\).  Hence \(x^L\GameLe x\).

  It remains to exclude the reverse comparison.  If \(x\GameLe x^L\), its
  left-option clause gives \(\neg(x^L\GameLe x^L)\), contradicting
  reflexivity from Theorem~\ref{thm:game-order-calculus}.  Thus
  \(x^L\GameLt x\) by Definition~\ref{def:game-order}.  The mirrored argument
  proves \(x\GameLt x^R\) for every right option \(x^R\).
\end{proof}

\begin{theorem}[Totality and trichotomy for surreal games]
  \label{thm:surreal-totality}
  \lean{IsSurreal.le_of_not_le,IsSurreal.totality,IsSurreal.trichotomy}
  \leanok
  If \(x\) and \(y\) are surreal, then
  \[
    x\GameLe y\ \vee\ y\GameLe x,
  \]
  and one of the three alternatives \(x\GameLt y\), \(x\GameEq y\), or
  \(y\GameLt x\) holds.
\end{theorem}
\begin{proof}
  \leanok
  \uses{def:game-order,def:is-surreal,thm:game-order-calculus,thm:between-options}
  Suppose \(\neg(x\GameLe y)\).  Negating the two universal clauses in
  Definition~\ref{def:game-order} produces one of two witnesses: either there
  is a left option \(x^L\) such that \(y\GameLe x^L\), or there is a right
  option \(y^R\) such that \(y^R\GameLe x\).  In the first case,
  Theorem~\ref{thm:between-options} gives \(x^L\GameLe x\); in the second it
  gives \(y\GameLe y^R\).  Transitivity from
  Theorem~\ref{thm:game-order-calculus} therefore yields \(y\GameLe x\) in
  either case.  Hence
  \[
    \neg(x\GameLe y)\quad\Longrightarrow\quad y\GameLe x,
  \]
  which proves totality.

  For trichotomy, choose a direction supplied by totality and then distinguish
  whether the reverse weak comparison also holds.  If both hold, then
  \(x\GameEq y\); if exactly one holds, then the corresponding comparison is
  strict.  These are precisely the definitions in
  Definition~\ref{def:game-order}.
\end{proof}

\begin{definition}[Surreal representatives]
  \label{def:surreal-representative}
  \lean{Surreal,Surreal.sr_zero,Surreal.sr_one,Surreal.leftOption,Surreal.rightOption}
  \leanok
  \uses{def:game,def:is-surreal}
  A surreal representative is a game together with a proof that it is surreal.  The surreal
  games \(0\) and \(1\) therefore give canonical representatives.  If \(x\) is a surreal
  representative, then each left or right option of its underlying game, equipped with the
  inherited surrealness assertion from Definition~\ref{def:is-surreal}, is again a surreal
  representative.  Equivalent game trees are not yet identified.
\end{definition}

\begin{definition}[Pair induction for surreal representatives]
  \label{def:surreal-pair-induction}
  \lean{Surreal.BiSurreal,Surreal.U,Surreal.wf_U,Surreal.U_left₁,Surreal.U_left₂,Surreal.U_right₁,Surreal.U_right₂}
  \leanok
  \uses{def:birthday,lem:birthday-options,def:surreal-representative}
  Order pairs \((a,b)\) of surreal representatives by
  \[
    (a',b')\mathrel{\vartriangleleft}(a,b)
      \quad\Longleftrightarrow\quad
    \bd(a')+\bd(b')<\bd(a)+\bd(b),
  \]
  where the birthday of a representative means the birthday of its underlying game.
  This relation is well founded.  Replacing either coordinate by any left or right option
  produces a smaller pair, by Lemma~\ref{lem:birthday-options}.
\end{definition}
